\section{Conditional recourse}\label{sec:recourse}
The condition numbers of Section~\ref{sec:growth} are built from the
coordinate curvatures $L_i$ of $F$. Stiff convex terms, which are easy to
optimize, can make them large, and so can a large residual problem whose
nonconvexity is combinatorial rather than metric. This section minimizes
exactly over part of the variables for fixed values of the others; we call
this \emph{recourse}, as for the second stage of a two-stage program,
although nothing here is random. The coordinates split into a
\emph{retained} set $\mathcal K$ and a \emph{recourse} set
$\mathcal R=[n]\setminus\mathcal K$, with $v=x_{\mathcal K}$ and
$y=x_{\mathcal R}$. For given $v$, $y$ is chosen optimally, and CT
(Algorithm~\ref{alg:ct}) runs on the value function $V$
of~\eqref{eq:valuefn}.

Section~\ref{sec:valuefn} shows that $V$ inherits coordinate curvature and
growth from $F$, so that Sections~\ref{sec:grids}--\ref{sec:exact} apply to
$V$ (Theorem~\ref{thm:valuefn}). Section~\ref{sec:convexrecourse} eliminates
private convex quadratic blocks and certifies the curvature of their value
factors (Theorem~\ref{thm:cv}). Section~\ref{sec:cuts} eliminates residual
problems that are concave in each coordinate and submodular after sign
changes, whose conditional values are minimum cuts
(Theorem~\ref{thm:cr-oracle}). These are the main results. Side results are
Section~\ref{sec:local}, which shows that a correction charged to one bag is
valid with exact recourse (Theorem~\ref{thm:cr-filter}) but not with grid
min-marginals (Proposition~\ref{prop:star}), and the extension of the cut
oracle to concave--convex residuals (Remark~\ref{rem:cr-mixed});
Section~\ref{sec:cuts} uses only
Lemma~\ref{lem:cr-cell} and Theorem~\ref{thm:cr-filter} from
Section~\ref{sec:local}. Section~\ref{sec:balanced} does not use recourse:
for balanced quadratics the grid problem of CT is itself a minimum cut, so
no tree decomposition is needed (Theorem~\ref{thm:balanced}).

\subsection{Value functions inherit curvature and growth}\label{sec:valuefn}
For the retained set $\mathcal K$ and the recourse set
$\mathcal R=[n]\setminus\mathcal K$ define
\begin{equation}\label{eq:valuefn}
V(v)=\min_{y\in X_{\mathcal R}}F(v,y),\qquad v\in\bar X_{\mathcal K},
\end{equation}
where $\bar X_{\mathcal K}=\prod_{i\in\mathcal K}[\ell_i,u_i]$ and the
feasible set $X_{\mathcal R}$ of the recourse problem is compact and does
not depend on $v$; this is the only structural requirement. For
model~\eqref{eq:model}, $X_{\mathcal R}=\prod_{i\in\mathcal R}X_i$; in
Section~\ref{sec:convexrecourse} it is a product of polytopes. We have
$\min_{X_{\mathcal K}}V=\OPT$, and a retained point $v\in X_{\mathcal K}$
together with a minimizer $y(v)$ of the recourse problem is a feasible point
with value $V(v)$.

\begin{lemma}[Partial minimization]\label{lem:valuefunction}
Let $V$ be defined by~\eqref{eq:valuefn}.
\begin{enumerate}[label=(\alph*)]
\item Let $i\in\mathcal K$ and $L_i\ge0$. If
$t\mapsto F(v+te_i,y)-\frac{L_i}{2}t^2$ is concave on
$\{t:v+te_i\in\bar X_{\mathcal K}\}$ for every $v\in\bar X_{\mathcal K}$
and every $y\in X_{\mathcal R}$, then so is
$t\mapsto V(v+te_i)-\frac{L_i}{2}t^2$ for every $v\in\bar X_{\mathcal K}$.
In particular, upper coordinate curvatures of $F$ for the coordinates in
$\mathcal K$ are upper coordinate curvatures of $V$ on $\bar X_{\mathcal K}$.
\item If $F$ has point growth with constant $g$ at $(v^*,y^*)$, then $V$
has point growth with constant $g$ at $v^*$ on $X_{\mathcal K}$.
\end{enumerate}
\end{lemma}

\begin{proof}
(a) For fixed $v$, the function $t\mapsto V(v+te_i)-\frac{L_i}{2}t^2$ is the
pointwise minimum over $y\in X_{\mathcal R}$ of the concave functions
$t\mapsto F(v+te_i,y)-\frac{L_i}{2}t^2$. It is finite, and it is concave
because its hypograph is the intersection of their hypographs.
(b) Since $\OPT\le V(v^*)\le F(v^*,y^*)=\OPT$, we have $V(v^*)=\OPT$. For
$v\in X_{\mathcal K}$ let $y(v)$ attain the minimum. Then
$V(v)-V(v^*)=F(v,y(v))-\OPT\ge g(\norm{v-v^*}^2+\norm{y(v)-y^*}^2)\ge
g\norm{v-v^*}^2$.
\end{proof}

Part~(a) is the classical stability of semiconcavity under infima
\cite{CannarsaSinestrari2004}. Lemma~\ref{lem:valuefunction}(a) explains why
the corrected-grid bound survives partial minimization, and
Lemma~\ref{lem:valuefunction}(b) why the growth analysis does. The curvature
of $V$ can be much smaller than that of $F$. If $F(v,y)=M(y-v)^2$ with
$v,y\in[0,1]$, the diagonal curvature in $v$ is $2M$, whereas $V\equiv0$.
Section~\ref{sec:convexrecourse} evaluates $V$ exactly and certifies
curvature bounds that can cancel such stiff convex terms
(Theorem~\ref{thm:cv}), though not always (Proposition~\ref{prop:cv-limit});
Section~\ref{sec:cuts} evaluates $V$ exactly when the eliminated problem is
combinatorial.

The factor structure survives the elimination as follows. Suppose the
recourse coordinates form disjoint \emph{private blocks}
$y^{(1)},\dots,y^{(r)}$ with domains $Y_1,\dots,Y_r$, so that
$X_{\mathcal R}=\prod_tY_t$, and
\begin{equation}\label{eq:privateblocks}
F(v,y)=F_0(v)+\sum_{t=1}^r\phi_t\bigl(v_{E_t},y^{(t)}\bigr),
\end{equation}
where $F_0$ is a sum of factors in $v$ and each block interacts with the
retained coordinates only through its \emph{scope} $E_t\subseteq\mathcal K$.
Then
\[
V(v)=F_0(v)+\sum_{t=1}^r\psi_t(v_{E_t}),\qquad
\psi_t(w)=\min_{y^{(t)}\in Y_t}\phi_t(w,y^{(t)}),
\]
is again a sum of factors, the \emph{value factors} $\psi_t$, with scopes
$E_t$. Any tree decomposition of the scopes of $F_0$ and the sets $E_t$
serves for $V$. Its width is the width after elimination of the blocks,
which can differ from the width of the original decomposition.

\begin{theorem}[Grids over value functions]\label{thm:valuefn}
Suppose $V$ has the form above, with a tree decomposition of the scopes of
$F_0$ and the sets $E_t$ of maximum bag size $p$, and with rational upper
coordinate curvatures $L_i^V$ $(i\in\mathcal K)$ of $V$ whose encoding
lengths are bounded by a polynomial in $I$. Suppose that, for every rational
point $w$, the factors of $F_0$ can be evaluated exactly, and an oracle
returns $\psi_t(w)$ and a minimizer $y^{(t)}(w)$ exactly for every block $t$,
both in time bounded by a polynomial in $I$ and the encoding length of $w$.
Put $L^V=\max_iL_i^V$.
\begin{enumerate}[label=(\alph*)]
\item Theorem~\ref{thm:certificate} holds for $V$, and the incumbent of CT
applied to $V$ lifts to a feasible point $(\hat v,y(\hat v))$ of value
$V(\hat v)$.
\item If $V$ has point growth with constant $g$ (by
Lemma~\ref{lem:valuefunction}(b), this holds if $F$ has), then
Theorem~\ref{thm:approx} holds for $V$ with $\kappa_V=\kappa(L^V,g)$ in
place of $\bar\kappa$ and with the bound
$f(p,\kappa_V)\,\kappa_V^{C}(I+q+1)^{C}$ in place of
$f(p,\bar\kappa)(I+q+1)^5$, where $C$ depends only on the polynomials in
the hypotheses.
\item If $F$ is a rational quadratic on the mixed box $X$
of~\eqref{eq:model}, Theorem~\ref{thm:exact} holds for EX
(Algorithm~\ref{alg:ex}) with CT applied to $V$ and with REC
(Algorithm~\ref{alg:rec}) and the acceptance test applied to $F$ at the
lifted incumbent. If $F$ has point
growth with constant $g$, the bound becomes
$f(p,\kappa_V)\,\kappa_V^{C_1}(I+1)^{C_1}$, where $C_1$ depends only on the
polynomials in the hypotheses.
\end{enumerate}
\end{theorem}

\begin{proof}
(a) Sections~\ref{sec:grids} and~\ref{sec:growth} use the objective only
through Definition~\ref{def:curvature}, exact factor values at grid points,
and, for the running time, growth. The $L_i^V$ satisfy
Definition~\ref{def:curvature} for $V$, the oracle gives exact values and
$y(\hat v)$, and $\min_{X_{\mathcal K}}V=\OPT$.

(b), (c) The accounting of bit lengths, which replaces the common
denominators of Theorem~\ref{thm:approx} by the bounds of the oracle, and of
the accuracy needed for exact output is in
Appendix~\ref{app:recourse-convex}.
\end{proof}

The theorem charges the work of the recourse oracle to the polynomial
factor, so it is useful when the recourse problems are easy and their
elimination does not increase the width. Sections~\ref{sec:convexrecourse}
and~\ref{sec:cuts} give classes in which the recourse problems are solved
exactly and with certificates.

\subsection{Bag-local corrections with exact or certified recourse}\label{sec:local}
The correction of Section~\ref{sec:grids} is global: every coordinate is
rounded, and the allowance is the sum of all coordinate corrections. A
bag-local allowance is valid when the coordinates outside the bag are
minimized exactly or with a certified error.

\paragraph{Bag cells.}
Let $B$ be a nonempty set of coordinates, for example a bag or the core of
Section~\ref{sec:cuts}. We write $V_B$ for the value
function~\eqref{eq:valuefn} with $\mathcal K=B$, and
$\bar X_B=\prod_{i\in B}[\ell_i,u_i]$. Let $L_i$ $(i\in B)$ be upper
coordinate curvatures of $F$ for the coordinates in $B$. A \emph{bag cell}
is a product $C=\prod_{i\in B}[a_i,b_i]$ with $a_i\le b_i$ and
$a_i,b_i\in X_i$. Its \emph{corners} are the points $v$ with
$v_i\in\{a_i,b_i\}$; they lie in $X_B$. With the effective width $w$ of
Section~\ref{sec:grids}, put
\[
e_B(C)=\sum_{i\in B}\frac{L_i}8\,w([a_i,b_i])^2 .
\]
If every side of $C$ is at most $h$, then $e_B(C)\le\frac{h^2}8\sum_{i\in B}L_i$.

\begin{lemma}[Bag-cell interpolation]\label{lem:cr-cell}
Let $C$ be a bag cell. Every $x\in X$ with $x_B\in C$ satisfies
\[
F(x)\ \ge\ \min_{v\text{ a corner of }C}V_B(v)-e_B(C).
\]
\end{lemma}
\begin{proof}
By Lemma~\ref{lem:valuefunction}(a), the $L_i$ satisfy
Definition~\ref{def:curvature} for $V_B$ on $\bar X_B$. The proof of
Proposition~\ref{prop:cellwise}(a) uses only Definition~\ref{def:curvature};
applied to $V_B$, to the subbox spanned by $C$ and to the grid of its
corners, it gives $V_B(x_B)\ge\min_{\text{corners}}V_B-e_B(C)$. Finally
$F(x)\ge V_B(x_B)$ because $x_{[n]\setminus B}\in X_{[n]\setminus B}$.
\end{proof}

No coordinate outside $B$ is rounded, so the error $e_B(C)$ does not grow
with $n$. The price is that $V_B$ is a global optimization problem over the
original outside domain.

\begin{theorem}[Conditional-recourse filter]\label{thm:cr-filter}
Let $\mathcal Q$ be a finite family of bag cells,
$e\ge\max_{C\in\mathcal Q}e_B(C)$ and $\varepsilon_{\mathrm{or}}\ge0$. Suppose that for every
corner $v$ of every cell of $\mathcal Q$ an oracle supplies a number
$\underline V(v)$ and a point $x^v\in X$ with $x^v_B=v$ such that
\begin{equation}\label{eq:cr-oracle}
\underline V(v)\le V_B(v),\qquad F(x^v)\le\underline V(v)+\varepsilon_{\mathrm{or}} .
\end{equation}
Let $U$ be the value of a feasible point with $U\le\min_vF(x^v)$, the
minimum taken over all supplied completions. Retain $C\in\mathcal Q$ if and
only if $\min_{v\text{ corner of }C}\underline V(v)-e\le U$. Then:
\begin{enumerate}
\item[(a)] every $C\in\mathcal Q$ that contains $x^*_B$ for some global
minimizer $x^*$ is retained;
\item[(b)] if some $C\in\mathcal Q$ contains $x^*_B$ for a global minimizer
$x^*$, then $U\le\OPT+e+\varepsilon_{\mathrm{or}}$, and every retained cell has a corner $v$
with $V_B(v)\le\OPT+2e+2\varepsilon_{\mathrm{or}}$;
\item[(c)] every $x\in X$ with $x_B\in\bigcup\mathcal Q$ satisfies
\[
F(x)\ge\lambda:=\min_{C\in\mathcal Q}\Bigl(\min_{\text{corners of }C}\underline V-e\Bigr),
\]
and $U-\lambda\le e+\varepsilon_{\mathrm{or}}$.
\end{enumerate}
\end{theorem}

The proof is in Appendix~\ref{app:recourse-convex}.

Appendix~\ref{app:recourse-convex} applies the theorem level by level to
multilevel schemes and shows that filtering alone needs only values whose
error varies by at most $\varepsilon_{\mathrm{or}}$ over the corners
(Proposition~\ref{prop:cr-osc}). Grid min-marginals, which minimize over
the outside coordinates on the grid instead of their domain, do not suffice:
a constant correction charged to one bag then keeps a cell containing the bag
projection of the minimizer only if the constant is at least a quantity that
grows with the number of outside coordinates, already for bag size two and
$\kappa<11$
(Proposition~\ref{prop:star} in Appendix~\ref{app:recourse-convex}).

\subsection{Certified convex recourse}\label{sec:convexrecourse}

Large positive curvature often sits in variables that enter the objective
convexly once the other variables are fixed. This subsection eliminates such
variables exactly, as convex value factors, and shows that CT then depends
only on the coordinate curvature and the growth of the reduced objective.
Proposition~\ref{prop:vf-curv} covers three cases (no certificate, a
globally affine response, a piecewise-affine response), and
Proposition~\ref{prop:cv-limit} shows the limits of this approach.

\begin{definition}[Convex recourse model]\label{def:cv-model}
The coordinates split into retained coordinates $\mathcal K$, with
$v=x_{\mathcal K}\in X_{\mathcal K}=\prod_{i\in\mathcal K}X_i$ and $X_i$ as
in~\eqref{eq:model}, and $r$ private blocks $y^{(t)}\in Y_t$,
$t=1,\ldots,r$, where $Y_t=\{y\in\R^{n_t}:G_ty\le b_t\}$ is a nonempty
bounded rational polytope. Private variables are continuous. Block $t$ has
an ordered scope $E_t\subseteq\mathcal K$ of $k_t$ retained coordinates; for
$i\in E_t$ let $j(t,i)$ be its position in $E_t$, and let
$\Pi_t=\prod_{i\in E_t}[\ell_i,u_i]$ be the \emph{scope box}. The objective
has the form~\eqref{eq:privateblocks},
\begin{equation}\label{eq:cv-model}
F(v,y)=F_0(v)+\sum_{t=1}^r\phi_t\bigl(v_{E_t},y^{(t)}\bigr),\qquad
\phi_t(w,y)=\tfrac12y^{\T}C_ty+y^{\T}\Gamma_tw+c_t^{\T}y,
\end{equation}
with rational data, $C_t\succeq0$, and a rational quadratic
$F_0(v)=\sum_{a\in\mathcal A}f_a(v_{S_a})$ with Hessian $H_0$. All terms
that involve retained variables only are placed in $F_0$; thus $(H_0)_{ii}$
is the diagonal entry of the full Hessian of $F$ at $v_i$. A tree
decomposition of the retained coordinates is supplied, with maximum bag size
$p$, such that every $S_a$ and every $E_t$ lies in a bag.
\end{definition}

The full Hessian of $F$ need not be positive semidefinite, and $C_t$ may be
singular. The value factors and the reduced objective are
\begin{equation}\label{eq:cv-value}
\psi_t(w)=\min_{y\in Y_t}\phi_t(w,y)\quad(w\in\R^{k_t}),\qquad
V(v)=F_0(v)+\sum_t\psi_t(v_{E_t}).
\end{equation}
Since the blocks are independent once $v$ is fixed, $V$ is the value
function~\eqref{eq:valuefn} with $X_{\mathcal R}=\prod_tY_t$, and
$\min_{X_{\mathcal K}}V=\OPT$. Every $v$ lifts to $(v,y(v))$ with
$y^{(t)}(v)\in\argmin_{Y_t}\phi_t(v_{E_t},\cdot)$ and $F(v,y(v))=V(v)$.

\paragraph{Leaf certificates and the curvature of a value factor.}
Fix one block, write $k=k_t$, and drop the index $t$:
$\phi(w,y)=\frac12y^{\T}Cy+y^{\T}\Gamma w+c^{\T}y$ for $w\in\R^k$,
$Y=\{y:Gy\le b\}$ with $m$ rows, and $\psi(w)=\min_{y\in Y}\phi(w,y)$.

\begin{definition}[Leaf certificate]\label{def:leaf}
Let $\Pi'\subseteq\R^k$ be a polytope with nonempty interior. A
\emph{certificate on $\Pi'$} consists of affine maps
$\bar y(w)=\bar y_0+Bw$ and $\lambda(w)=\lambda_0+\Lambda w\in\R^m$ such
that, for all $w\in\Pi'$,
\begin{align}
&G\bar y(w)\le b,\qquad \lambda(w)\ge0,\label{eq:leaf-feas}\\
&C\bar y(w)+\Gamma w+c+G^{\T}\lambda(w)=0,\label{eq:leaf-stat}\\
&\lambda_s(w)\bigl(b_s-G_s\bar y(w)\bigr)=0\qquad(s=1,\ldots,m).\label{eq:leaf-comp}
\end{align}
A \emph{certificate on a rational box $\Pi\subseteq\R^k$} is a rational
binary space partition of $\Pi$, in which each internal node splits its
polytope by one rational hyperplane into two closed halves with nonempty
interiors, together with a certificate on each of its leaves
$\Pi'_1,\ldots,\Pi'_{N_\Pi}$. For a private box $Y=[\ell^Y,u^Y]$ the rows are
$y\le u^Y$ and $-y\le-\ell^Y$, and $\lambda$ splits into upper and lower
bound multipliers.
\end{definition}

A certificate is checked exactly in time polynomial in its encoding length.
The leaves of a binary space partition cover its root, each leaf is
described by the inequalities of $\Pi$ and of the halfspaces on its root
path, and one linear program per node checks that both halves have
nonempty interiors. Stationarity~\eqref{eq:leaf-stat} is an identity between
affine maps, and complementarity~\eqref{eq:leaf-comp} says that a quadratic
polynomial vanishes on a set with nonempty interior, that is, all its
coefficients vanish. By linear programming duality, an affine inequality
$\alpha_0+\alpha^{\T}w\ge0$ holds on a nonempty polytope $\{w:Aw\le a\}$ if
and only if some $\pi\ge0$ has $A^{\T}\pi=-\alpha$ and
$a^{\T}\pi\le\alpha_0$; such rational multipliers are part of the
certificate.

\begin{lemma}[Certificate identity and piece Hessian]\label{lem:leaf}
Let $(\bar y,\lambda)$ be a certificate on $\Pi'$. Then:
\begin{enumerate}
\item[(a)] For $w\in\Pi'$ and $y\in\R^{n_t}$,
\begin{equation}\label{eq:kkt-identity}
\phi(w,y)-\phi(w,\bar y(w))=\tfrac12(y-\bar y(w))^{\T}C(y-\bar y(w))
+\lambda(w)^{\T}(b-Gy).
\end{equation}
The right side is nonnegative for $y\in Y$. Hence $\psi=q_{\Pi'}$ on $\Pi'$,
where $q_{\Pi'}(w)=\phi(w,\bar y(w))$.
\item[(b)] Let $\mathcal E=\{s:\lambda_s\not\equiv0\}$. Then $G_sB=0$ for
$s\in\mathcal E$, and
\[
\nabla^2q_{\Pi'}=B^{\T}CB+B^{\T}\Gamma+\Gamma^{\T}B=-B^{\T}CB\preceq0.
\]
\end{enumerate}
\end{lemma}

\begin{proof}
(a) The function $\phi(w,\cdot)$ is quadratic with Hessian $C$, and
by~\eqref{eq:leaf-stat} its gradient at $\bar y=\bar y(w)$ is
$-G^{\T}\lambda$. Exact expansion gives
\[
\phi(w,y)-\phi(w,\bar y)=\lambda^{\T}(G\bar y-Gy)+\tfrac12(y-\bar y)^{\T}C(y-\bar y),
\]
and $\lambda^{\T}(G\bar y-Gy)=\lambda^{\T}(b-Gy)-\lambda^{\T}(b-G\bar y)
=\lambda^{\T}(b-Gy)$ by~\eqref{eq:leaf-comp}. For $y\in Y$ both terms on the
right of~\eqref{eq:kkt-identity} are nonnegative, since $C\succeq0$ and
$\lambda(w)\ge0$. Equality holds at $y=\bar y(w)\in Y$.

(b) For each $s$, the polynomial $\lambda_s(w)(b_s-G_s\bar y_0-G_sBw)$
vanishes on $\Pi'$, which has nonempty interior, so it is the zero
polynomial. The polynomial ring is an integral domain; for $s\in\mathcal E$
the first factor is a nonzero polynomial, so the second is zero and
$G_sB=0$. The $w$-linear part of~\eqref{eq:leaf-stat} reads
$CB+\Gamma+G^{\T}\Lambda=0$, and row $s$ of $\Lambda$ vanishes for
$s\notin\mathcal E$. Hence
\[
B^{\T}(CB+\Gamma)=-\sum_{s\in\mathcal E}(G_sB)^{\T}\Lambda_s=0,
\]
so $B^{\T}\Gamma=-B^{\T}CB$, which is symmetric; therefore also
$\Gamma^{\T}B=-B^{\T}CB$. Substituting $y=\bar y_0+Bw$ in $\phi$ gives
$\nabla^2q_{\Pi'}=B^{\T}CB+B^{\T}\Gamma+\Gamma^{\T}B=-B^{\T}CB$.
\end{proof}

By~\eqref{eq:cv-energy} in Appendix~\ref{app:recourse-convex},
$(B^{\T}CB)_{jj}$ is the energy that the private response spends to follow a
unit change of $w_j$.

\begin{proposition}[Coordinate curvature of a value factor]\label{prop:vf-curv}
\begin{enumerate}
\item[(a)] $\psi$ is finite and concave on $\R^k$.
\item[(b)] Let a certificate on a box $\Pi$ be given, with linear part
$B_{\Pi'}$ of the response on the leaf $\Pi'$, and put
\[
\sigma_j=\min_{\Pi'}\,(B_{\Pi'}^{\T}CB_{\Pi'})_{jj}\ \ge0,
\]
the minimum over all leaves. For every $w\in\Pi$ and every interval $J$
with $w+Je_j\subseteq\Pi$, the function $t\mapsto\psi(w+te_j)+\sigma_jt^2/2$
is concave on $J$.
\item[(c)] If a number $\sigma'$ has the property in (b), then
$\sigma'\le\sigma_j$. In particular $\sigma_j$ is the same for every
certificate on $\Pi$.
\end{enumerate}
\end{proposition}

\begin{proof}
(a) For fixed $y$, $\phi(w,y)$ is affine in $w$. Thus $\psi$ is a pointwise
minimum of affine functions over the compact nonempty set $Y$, hence finite
and concave.

(b) On each leaf, $\psi=q_{\Pi'}$ has second derivative
$-(B_{\Pi'}^{\T}CB_{\Pi'})_{jj}\le-\sigma_j$ along $e_j$
(Lemma~\ref{lem:leaf}), and at the finitely many points where a segment
changes leaf, the one-sided derivatives of the concave function $\psi$ can
only decrease. Appendix~\ref{app:recourse-convex} gives the details.

(c) Take any leaf $\Pi'$ and an interior point $w'$ of $\Pi'$. For small
$|t|$, $\psi(w'+te_j)=q_{\Pi'}(w'+te_j)$, whose second derivative is
$-(B_{\Pi'}^{\T}CB_{\Pi'})_{jj}$. Concavity of $\psi(w'+te_j)+\sigma't^2/2$
forces $\sigma'\le(B_{\Pi'}^{\T}CB_{\Pi'})_{jj}$. Minimizing over the leaves
gives $\sigma'\le\sigma_j$. Applying this to two certificates shows that
each $\sigma_j$ is at most the other.
\end{proof}

Part~(a) alone certifies $\sigma_j=0$ for every block; this is the direct
bound. Part~(c) shows that a piecewise certificate loses nothing: the
cancellation it certifies is the largest one possible for coordinatewise
semiconcavity of this factor. For a private box, a certificate on every
rational box always exists, so this largest cancellation is always
certifiable (Proposition~\ref{prop:cert-exist} in
Appendix~\ref{app:recourse-convex}). The construction can be exponential; no
bound on the number of leaves is claimed.

\begin{theorem}[Certified convex recourse]\label{thm:cv}
In the model of Definition~\ref{def:cv-model}, let each block carry a
certificate on its scope box $\Pi_t$, or none. Put $\sigma_t=0$ for a block
without certificate and otherwise let $\sigma_t$ be the vector of
Proposition~\ref{prop:vf-curv}(b). Define
\begin{equation}\label{eq:cv-L}
L_i=(H_0)_{ii}-\sum_{t:\,i\in E_t}\sigma_{t,j(t,i)},\qquad
L_i^+=\max\{L_i,0\},\qquad L^+=\max_iL_i^+,
\end{equation}
and let $I$ be the total encoding length of data, decomposition and
certificates.
\begin{enumerate}
\item[(i)] For each $i\in\mathcal K$, the function $v\mapsto V(v)-L_iv_i^2/2$
is concave on every segment of $\bar X_{\mathcal K}$ parallel to $e_i$. All
certificates are checked in $\poly(I)$ time.
\item[(ii)] If $L^+=0$, that is, $L_i\le0$ for all $i$, endpoint dynamic
programming with two nodes per coordinate returns an exact minimizer of $F$
in $2^{O(p)}\poly(I)$ bit operations. No growth or uniqueness is needed.
\item[(iii)] Suppose that $L^+>0$ and that $V$ has point growth with
constant $g$, as it has if $F$ has (Lemma~\ref{lem:valuefunction}(b)).
With $\kappa_V=\kappa(L^+,g)$, CT applied to $V$ with the curvatures $L_i^+$
returns a feasible original point with a certified gap at most $2^{-q}$ in
$f(p,\kappa_V)\kappa_V^C(I+q+1)^C$ bit operations. If $F$ has point growth
with constant $g$, an exact rational minimizer is found in
$f(p,\kappa_V)\kappa_V^{C_1}(I+1)^{C_1}$ bit operations. Here $f$ is the
function of Theorem~\ref{thm:approx}, $C$ and $C_1$ are absolute constants,
and neither $g$ nor $\kappa_V$ is an input. Every certified lower bound and
every returned exact minimizer is valid without growth or uniqueness.
Without point growth of $F$, the exact-output procedure has no termination
guarantee; if every $Y_t$ is a box, Theorem~\ref{thm:valuefn}(c) applies
instead, and EX stops on every instance.
\item[(iv)] $L_i\le(H_0)_{ii}$ for every $i$. Hence, if $F$ has point
growth with constant $g$ and $L'$ is a common upper coordinate curvature of
$F$ for all coordinates, retained and private, then
$\kappa(L^+,g)\le\kappa(L',g)$: the parameter $\kappa_V$ of (iii) never
exceeds the Euclidean condition number of $F$ with the private variables
kept as ordinary coordinates.
\end{enumerate}
\end{theorem}

The proof is in Appendix~\ref{app:cv-proof}.

\begin{remark}[Three instances and prior work]\label{rem:cv-instances}
Without certificates, $L_i=(H_0)_{ii}$: changing active sets are handled
exactly, but direct stiffness is not cancelled. A one-leaf certificate is a
globally affine response; then $\psi_t$ is an explicit quadratic with
Hessian $-B_t^{\T}C_tB_t$ that can be merged into $F_0$ (for $M(y-v)^2$ on
$[0,1]^2$ it gives $\sigma=2M$ and $L_v=0$). A certificate with several
leaves is a piecewise-affine response and certifies the largest possible
$\sigma_j$. The partitions of different blocks are never overlaid; their
common refinement can have exponentially many cells. Parametric-QP critical
regions are classical \cite{BemporadEtAl2002}; the contribution here is
their curvature accounting inside CT. Khajavirad \cite{Khajavirad2026PolyBox}
eliminates the components formed by variables with positive diagonal into
value functions of their neighbors; the neighbors have nonpositive diagonal,
can be taken binary, and their values are enumerated. Here the neighbors
$v_{E_t}$ are gridded continuous or integer coordinates, and what is
certified is the coordinate curvature of each value factor.
\end{remark}

Appendix~\ref{app:recourse-convex} derives an affine condensation
(Corollary~\ref{cor:cv-affine}), the bound
$\kappa_V\le\max\{1,C_0\}\cdot\max\{1,\nu/g\}$ when $L^+\le C_0\nu$ with
$\nu=\max\{0,-\lambda_{\min}(\nabla^2F)\}$ (Corollary~\ref{cor:cv-nu}), a
block whose certified curvature is $8$ while its direct curvature is $2M+8$
(Example~\ref{ex:cv-fm}), and a ladder with $m$ negative eigenvalues whose
certified ratio $L^+/g$ is at most $123$ for all $m$ and $M$
(Proposition~\ref{prop:cv-ladder}). For a block whose domain is a box, one
exact convex quadratic program and one linear program decide whether some
affine map is an optimal response on the whole scope box, and if so return
it with a one-leaf certificate (Theorem~\ref{thm:cv-recog} in
Appendix~\ref{app:recourse-convex}).

\paragraph{Limits of global recourse curvature.}
Call a retained set $\mathcal K\subseteq[n]$ \emph{admissible} if the
principal Hessian block of the coordinates in $\mathcal R=[n]\setminus\mathcal K$
is positive semidefinite, so that they form convex private blocks. For
admissible $\mathcal K$ the value function $V_{\mathcal K}$ is intrinsic: it
does not depend on how the private part is grouped or certified. Let
$L_{\mathcal K}$ be the least $L$ for which Definition~\ref{def:curvature}
holds for $V_{\mathcal K}$ in every retained coordinate, and
$g_{\mathcal K}$ the largest growth constant of $V_{\mathcal K}$. Since every
certified $L_i$ is valid (Theorem~\ref{thm:cv}(i)), $L^+\ge L_{\mathcal K}$,
and every growth constant $g$ of $V_{\mathcal K}$, in particular every growth
constant of $F$ (Lemma~\ref{lem:valuefunction}(b)), satisfies
$g\le g_{\mathcal K}$; so no certificate yields a parameter below
$L_{\mathcal K}/g_{\mathcal K}$.

\begin{proposition}[Recourse cannot replace $L/g$ by $\nu/g$ in general]\label{prop:cv-limit}
For $M\ge1$ let $x\in[0,2]$, $y\in[1,3]$ and
\[
F_M(x,y)=M(x-y)^2-\tfrac18(x+y-2)^2+2(y-1).
\]
Then $(1,1)$ is the unique minimizer, $\OPT=0$,
$F_M(x,y)-\OPT\ge\frac18\bigl((x-1)^2+(y-1)^2\bigr)$, and $\nu=\frac12$; thus
$\nu/g\le4$ with $g=\frac18$. Nevertheless, for every admissible retained
set $\mathcal K$ and every positive diagonal rescaling of the retained
coordinates,
\[
L_{\mathcal K}/g_{\mathcal K}\ \ge\ (4M-\tfrac12)/3\ \ge\ M.
\]
\end{proposition}

The proof is in Appendix~\ref{app:recourse-convex}, which also locates the
obstruction in a clipped piece of the private response. Disjoint copies give
arbitrarily many variables with the same constants and bag size two. The
instance is trivial to solve; the proposition limits the parameter of
Theorem~\ref{thm:cv}, not the problem.

\subsection{Separately concave residuals and minimum cuts}\label{sec:cuts}

In this subsection $F(x)=\frac12x^{\T}Hx+b^{\T}x+c$ is a rational quadratic,
the retained set $\mathcal K$ is a \emph{core} of $k$ continuous
coordinates, and the recourse set $\mathcal R$ is called the
\emph{residual}. Assume
\begin{equation}\label{eq:cr-class}
\text{(R1)}\ H_{ii}\le0\ (i\in\mathcal R),\qquad
\text{(R2)}\ o_io_jH_{ij}\le0\ (i\ne j,\ i,j\in\mathcal R)\ \text{for some }
o\in\{-1,1\}^{\mathcal R}.
\end{equation}
Nothing is assumed about the entries of $H$ in core rows, and the residual
interaction graph may be dense. A residual coordinate may be continuous or
integer; integer bounds are rounded inward first, and an empty domain is
rejected. Signs $o$ as in (R2) are found, or shown not to exist, by the
following lemma applied to $H_{\mathcal R\mathcal R}$.

\begin{lemma}[Recognizing balance]\label{lem:balance}
For a symmetric $m\times m$ matrix $H$, signs $o\in\{-1,1\}^m$ with
$o_io_jH_{ij}\le0$ for all $i\ne j$ can be found, or shown not to exist, in
time linear in $m$ plus the number of nonzero off-diagonal entries.
\end{lemma}

\begin{proof}
The conditions are $o_i=o_j$ when $H_{ij}<0$ and $o_i=-o_j$ when
$H_{ij}>0$. Breadth-first search fixes one sign in each connected component
of the graph of nonzero off-diagonal entries, propagates the conditions
along a spanning tree and checks every other edge. If an edge violates its
condition, it closes a cycle with an odd number of positive entries, around
which the conditions are inconsistent \cite{Harary1953}; otherwise all
conditions hold.
\end{proof}

\begin{lemma}[Endpoint reduction]\label{lem:cr-endpoint}
Fix $v\in\bar X_{\mathcal K}$ and a box $\prod_{i\in\mathcal R}[\ell_i,u_i]$
whose bounds are integers for integer coordinates. Under (R1), the minimum
of $F(v,\cdot)$ over this continuous box is attained at a point with
$y_i\in\{\ell_i,u_i\}$ for all $i\in\mathcal R$, and it equals the minimum
over the mixed residual domain.
\end{lemma}

\begin{proof}
Start from a minimizer, which exists by compactness. For each residual
coordinate in turn, $F$ as a function of that coordinate alone is
$\frac12H_{ii}t^2$ plus an affine function of $t$, which is concave by
(R1); replace the coordinate by an endpoint with no larger value. The
result is a minimizer whose residual coordinates are endpoints. Its integer
coordinates are integral, so it lies in the mixed domain, which is
contained in the continuous box.
\end{proof}

For $i\in\mathcal R$ put $\alpha_i=\ell_i$, $\delta_i=u_i-\ell_i$ if $o_i=1$,
and $\alpha_i=u_i$, $\delta_i=-(u_i-\ell_i)$ if $o_i=-1$. For a
\emph{label} $z\in\{0,1\}^{\mathcal R}$ let $y(z)$ have coordinates
$y(z)_i=\alpha_i+\delta_iz_i$; as $z$ ranges over $\{0,1\}^{\mathcal R}$,
$y(z)$ ranges over all residual endpoint vectors. In a
network with source $\mathsf s$ and sink $\mathsf t$, $\operatorname{cap}(S)$
denotes the total capacity of the arcs leaving a node set $S$ that contains
$\mathsf s$ but not $\mathsf t$.

\begin{proposition}[Signed cut representation]\label{prop:cr-cut}
(a) Let $\Phi(z)=\Phi_0+\sum_i\mu_iz_i+\sum_{i<j}\omega_{ij}z_iz_j$ on
$\{0,1\}^m$ with $\omega_{ij}\le0$ for all $i<j$, and put
$\omega_{ji}=\omega_{ij}$ and $\rho_i=\mu_i+\frac12\sum_{j\ne i}\omega_{ij}$.
Build a network on $\{1,\dots,m,\mathsf s,\mathsf t\}$
with arcs $i\to\mathsf t$ of capacity $\max\{\rho_i,0\}$ and
$\mathsf s\to i$ of capacity $\max\{-\rho_i,0\}$ for each $i$, and arcs
$i\to j$ and $j\to i$ of capacity $-\omega_{ij}/2\ge0$ for each pair with
$\omega_{ij}\ne0$. With $S_z=\{\mathsf s\}\cup\{i:z_i=1\}$,
\begin{equation}\label{eq:cr-cutid}
\Phi(z)=\Phi_0+\sum_i\min\{0,\rho_i\}+\operatorname{cap}(S_z)
\qquad(z\in\{0,1\}^m).
\end{equation}
(b) Under~\eqref{eq:cr-class} and for a rational core point $v$, the
function $z\mapsto F(v,y(z))$ on $\{0,1\}^{\mathcal R}$ has the form of (a)
with rational coefficients obtained by exact evaluation of $F$, and
$\omega_{ij}=H_{ij}\delta_i\delta_j\le0$.
\end{proposition}

\begin{proof}
(a) For binary $z$,
$\omega_{ij}z_iz_j=\frac{\omega_{ij}}2(z_i+z_j)-\frac{\omega_{ij}}2\abs{z_i-z_j}$,
so $\Phi(z)=\Phi_0+\sum_i\rho_iz_i-\frac12\sum_{i<j}\omega_{ij}\abs{z_i-z_j}$. In the
cut $S_z$, an arc $i\to\mathsf t$ is cut exactly when $z_i=1$, an arc
$\mathsf s\to i$ exactly when $z_i=0$, and exactly one of the arcs
$i\to j$, $j\to i$ when $z_i\ne z_j$, none otherwise. Hence
$\operatorname{cap}(S_z)=\sum_i[\max\{\rho_i,0\}z_i+\max\{-\rho_i,0\}(1-z_i)]
-\frac12\sum_{i<j}\omega_{ij}\abs{z_i-z_j}$, and
$\rho_iz_i=\max\{\rho_i,0\}z_i+\max\{-\rho_i,0\}(1-z_i)+\min\{0,\rho_i\}$.

(b) Expand $F(v,\alpha+\delta\circ z)$ and use $z_i^2=z_i$. Terms without
residual coordinates are constant. The terms $\frac12H_{ii}y_i^2$, $b_iy_i$
and $H_{ij}v_jy_i$ with $j\in\mathcal K$ are affine in $z_i$, and
$H_{ij}y_iy_j$ with $i\ne j$ in $\mathcal R$ gives
$H_{ij}\delta_i\delta_jz_iz_j$ plus affine terms. By (R2),
$\omega_{ij}=H_{ij}o_io_j(u_i-\ell_i)(u_j-\ell_j)\le0$.
\end{proof}

Every node set that contains $\mathsf s$ but not $\mathsf t$ is $S_z$ for
exactly one $z$, so by~\eqref{eq:cr-cutid} minimizing $\Phi$ is a minimum
$\mathsf s$--$\mathsf t$ cut problem with nonnegative capacities.

\begin{theorem}[Certified cut recourse]\label{thm:cr-oracle}
Assume~\eqref{eq:cr-class}. For every rational core point $v$, one
maximum-flow computation returns a minimizer $y(v)$ of $F(v,\cdot)$ over
$X_{\mathcal R}$, which is a residual endpoint vector, the exact value
$V(v)$, and a certificate: a label $z$ with $y(v)=y(z)$ and a feasible flow
in the network of Proposition~\ref{prop:cr-cut} whose value equals
$\operatorname{cap}(S_z)$. The bit complexity is polynomial in the encoding
length of $F$ and $v$, with an absolute exponent. Given the network, the
certificate is checked with a number of rational operations linear in
$\abs{\mathcal R}$ plus the number of nonzero entries $H_{ij}$ with
$i\ne j$ in $\mathcal R$.
\end{theorem}

The proof is in Appendix~\ref{app:cr-proofs}.

\begin{remark}[Scope and prior work]\label{rem:cr-cut-prior}
Lemmas~\ref{lem:balance} and~\ref{lem:cr-endpoint} and
Proposition~\ref{prop:cr-cut} are classical (Section~\ref{sec:related}).
With an empty core, Theorem~\ref{thm:cr-oracle} is the polynomial
solvability of box QPs that are submodular after sign changes and have
nonpositive diagonal \cite{Padberg1989,BurerNatarajanWillemsen2026v3}. What
we use is the conditional form: the core and the core--residual couplings
are arbitrary, and each conditional value comes with a flow certificate.
Del Pia and Khajavirad also fix the coordinates with nonpositive diagonal at
endpoints; their polynomial classes \cite[Theorem~4]{DelPiaKhajavirad2026}
restrict the structure of these coordinates (logarithmic treewidth and
logarithmic interfaces to the others) and the rank of their coupling to the
others, whereas (R2) restricts only signs and allows a dense residual.
Condition (R1) is what reduces the residual to endpoint vectors. Without it,
and with continuous residual coordinates, the conditional problem is a box
QP that is submodular after sign changes, and the complexity of that class
is open \cite[footnote~2]{BurerNatarajanWillemsen2026v3};
Section~\ref{sec:balanced} grids such coordinates instead
(Corollary~\ref{cor:core}).
\end{remark}

\paragraph{Search over a continuous core.}
Let the core domain be $X_{\mathcal K}=[0,1]^k$ with $k\ge1$, and let the
residual domain $X_{\mathcal R}$ be any compact mixed box. Let $L\ge0$ be an
upper coordinate curvature of $F$ for every core coordinate (for a
quadratic, any $L\ge\max_{i\in\mathcal K}H_{ii}$); by
Lemma~\ref{lem:valuefunction}(a) it is one for $V$. An \emph{exact oracle}
returns, for every dyadic $v\in[0,1]^k$, the value $V(v)$ and a completion
$y(v)\in X_{\mathcal R}$ with $F(v,y(v))=V(v)$, together with a certificate.
Theorem~\ref{thm:cr-oracle} supplies one under~\eqref{eq:cr-class}, and
Remark~\ref{rem:cr-mixed} supplies one when the residual has a convex part.
Level $j$ of the search uses dyadic cells of side $h_j=2^{-j}$, the common
mesh of Section~\ref{sec:growth} with $s=1$, and
\[
e_j=\frac{kLh_j^2}8,\qquad
\beta_V(C)=\min\{V(v):v\text{ a corner of }C\}-e_j .
\]

\begin{algorithm}[CORE: core search]\label{alg:core}
Given an exact oracle, $L$, an accuracy $\varepsilon>0$ and a level limit
$J$, start with $\mathcal Q_0=\{[0,1]^k\}$ and $U=+\infty$. For
$j=0,1,\dots,J$:
\begin{enumerate}[label=(\roman*)]
\item query the oracle at every corner of a cell of $\mathcal Q_j$ that has
not been queried before, and let $U$ be the least value $V(v)$ of all
queried corners, attained by the incumbent $(v_U,y(v_U))$;
\item put $\lambda_j=\min_{C\in\mathcal Q_j}\beta_V(C)$ and retain
$\mathcal Q_j'=\{C\in\mathcal Q_j:\beta_V(C)\le U\}$;
\item if $U-\min\{\lambda_j,U\}\le\varepsilon$, stop; otherwise let
$\mathcal Q_{j+1}$ consist of the $2^k$ dyadic children of each cell of
$\mathcal Q_j'$.
\end{enumerate}
Return the incumbent, $U$ and the lower bound $\min\{\lambda_j,U\}$.
\end{algorithm}

\begin{theorem}[Core search]\label{thm:cr-search}
At every level $j$ of CORE:
\begin{enumerate}[label=(\alph*)]
\item every cell of $\mathcal Q_j$ that contains $x^*_{\mathcal K}$ for a
global minimizer $x^*$ is retained; in particular $\mathcal Q_j'\ne\emptyset$;
\item $\lambda_j\le\OPT\le U\le\lambda_j+e_j$;
\item every retained cell has a corner $v$ with $V(v)\le\OPT+2e_j$;
\item every $x\in X$ satisfies $F(x)\ge\min\{\lambda_j,U\}$, and this can be
verified from the list of cells, the certified queried values and the
retention decisions, without any optimization.
\end{enumerate}
\end{theorem}

The proof is in Appendix~\ref{app:cr-proofs}.

By Theorem~\ref{thm:cr-search}(b), CORE stops at the latest at the first
level $j$ with $e_j\le\varepsilon$, provided that $J$ is at least this
level. Its cost is controlled by the number of near-optimal corners.

\begin{proposition}[Query count under core growth]\label{prop:cr-growth}
Suppose $V(v)-\OPT\ge g\norm{v-v^*}^2$ for all $v\in[0,1]^k$, some $v^*$
and some $g>0$; by Lemma~\ref{lem:valuefunction}(b) this holds with the same
$g$ if $F$ has point growth with constant $g$. Let
$\kappa_{\mathcal K}=\kappa(L,g)=\max\{1,L/g\}$. Then levels $0,\dots,J$ of
CORE make at most
\[
2^k+8^kJ\bigl(1+\sqrt{k\kappa_{\mathcal K}}\bigr)^k
\]
oracle queries. The constant $g$ is not used by the algorithm.
\end{proposition}

The proof is in Appendix~\ref{app:cr-proofs}.

The bound is linear in the number of levels, and the residual enters only
the cost of each query. Without growth the count can be exponential in the
number of levels.

\begin{example}[Flat directions]\label{ex:cr-flat}
Take $k$ even, no residual, and
$V(v)=\sum_{t=1}^{k/2}(v_{2t-1}-v_{2t})^2$, so $L=2$ and $\OPT=0$. The corner $v=0$ gives $U=0$ at level $0$. Every product
of diagonal cells
$\prod_t\bigl([a_th_j,(a_t+1)h_j]\times[a_th_j,(a_t+1)h_j]\bigr)$ with
integers $0\le a_t<2^j$ has a
corner with $V=0$, hence is retained, and its parent is again such a
product. So level $j$ retains at least $2^{jk/2}$ cells, while
$U-\lambda_j=e_j$; stopping at accuracy $\varepsilon$ therefore takes
$\Omega(\varepsilon^{-k/4})$ queries for fixed $k$.
\end{example}

\paragraph{Exact output.}
Under~\eqref{eq:cr-class} and for rational data, $\OPT$ and every value $\min_{v\in[0,1]^k}F(v,y)$ at a
residual endpoint vector $y$ have denominators at most a constant
$\Omega_{\mathcal K}$ of bit length $O(I^2)$ (Lemma~\ref{lem:cr-height}).
Run CORE with the oracle of Theorem~\ref{thm:cr-oracle} and accuracy
$\varepsilon=\Omega_{\mathcal K}^{-2}/2$, and minimize $F(\cdot,y)$
exactly over the faces of the core cube, where $y$ is the residual part of
the incumbent. The result is a global minimizer that passes the acceptance
test of Proposition~\ref{prop:accept} with the denominator bound
$\Omega_{\mathcal K}$. Under core growth the total cost is
$8^k(1+\sqrt{k\kappa_{\mathcal K}})^k\poly(I)$ bit operations with an
absolute exponent, and $g$ is not an input (Theorem~\ref{thm:cr-exact} in
Appendix~\ref{app:recourse-cuts}).

\begin{remark}[Concave--convex residuals]\label{rem:cr-mixed}
The oracle extends to residuals with a convex part. Split
$\mathcal R=\mathcal R_-\cup\mathcal R_+$, require (R1) only on $\mathcal R_-$
and (R2) on all of $\mathcal R$, and let the coordinates in $\mathcal R_+$
be continuous with $H_{\mathcal R_+\mathcal R_+}\succeq0$. For a fixed core
point, the coordinates in $\mathcal R_-$ can be moved to endpoints without
increasing $F$, as in Lemma~\ref{lem:cr-endpoint}. For each set
$S\subseteq\mathcal R_-$ of coordinates with label $z_i=1$, what remains is a
convex box QP in $y_{\mathcal R_+}$, which is solved exactly in polynomial
time \cite{KozlovTarasovKhachiyan1980}. Its value $\Phi(S)$ is a submodular
function of $S$, because partial minimization preserves submodularity
\cite{Topkis1978}, and the conditional minimum is $\min_S\Phi(S)$.
Submodular minimization by the ellipsoid method
\cite{GrotschelLovaszSchrijver1988} and Edmonds' min--max theorem
\cite{Edmonds1970} then give the exact conditional value, a minimizer, and a
certificate: a convex combination of at most $\abs{\mathcal R_-}+1$ greedy
vectors of $\Phi$. All these ingredients are classical, and the same
composition underlies the mixed continuous--discrete methods cited in
Section~\ref{sec:related}.
Proposition~\ref{prop:cr-greedy} in Appendix~\ref{app:recourse-cuts} states
the result precisely. It supplies, for concave--convex residuals, the exact
oracle required by CORE (Theorem~\ref{thm:cr-search} and
Proposition~\ref{prop:cr-growth}), with a certificate that can be checked
without optimization (Proposition~\ref{prop:cr-greedy}(d)). For
$\mathcal R_+=\emptyset$ this is the class
\eqref{eq:cr-class}, where Theorem~\ref{thm:cr-oracle} needs only one
maximum flow.
\end{remark}

\subsection{Corrected grids without a tree decomposition}\label{sec:balanced}

CT (Algorithm~\ref{alg:ct}) uses its tree decomposition only to compute, on
the current product grid $G=\prod_iG_i$, the corrected minimum
$\beta=\min_GQ$, one minimizer and the min-marginals $m_i(v)$ of $Q=F-D$
(Lemma~\ref{lem:dp}). Any exact method for these quantities can replace the
dynamic program. For the following sign structure, minimum cuts are such a
method, whatever the interaction graph.

\begin{definition}[Balanced quadratic]\label{def:balanced}
A quadratic $F(x)=\frac12x^{\T}Hx+b^{\T}x+c$ is \emph{balanced} if there are
signs $o\in\{-1,1\}^n$ with $o_io_jH_{ij}\le0$ for all $i\ne j$.
\end{definition}

Equivalently, after the substitution $x_i\mapsto o_ix_i$ all off-diagonal
Hessian entries are nonpositive, so $F$ becomes submodular on $\R^n$.
Condition (R2) of Section~\ref{sec:cuts} says that $F$ becomes balanced once
the core coordinates are fixed, and Lemma~\ref{lem:balance} recognizes
balance in linear time.
No sign condition is imposed on the diagonal or on the linear terms. If
$H_{ii}\le0$ for all $i$, then $F$ satisfies~\eqref{eq:cr-class} with an
empty core, and Theorem~\ref{thm:cr-oracle} minimizes it exactly with one
maximum flow.

\begin{lemma}[Grid minimization by one minimum cut]\label{lem:chaincut}
Let $F$ be balanced with signs $o$, let $G_i\subset\Q$ be finite and
nonempty with $\abs{G_i}=k_i+1$, and let $\chi_i:G_i\to\Q$ be arbitrary. Put
$F_\chi(y)=F(y)+\sum_i\chi_i(y_i)$ for $y\in\prod_iG_i$. Then $\min F_\chi$
and a minimizer are obtained from one minimum $\mathsf s$--$\mathsf t$ cut in
a network with $\sum_ik_i+2$ nodes and at most
$3\sum_ik_i+2\sum_{i<j:\,H_{ij}\ne0}k_ik_j$ arcs, whose capacities are
computed from $H$, $b$, the grid nodes and the values of the $\chi_i$ by
polynomially many rational operations. A feasible flow of equal value
certifies the minimum. The same holds for each min-marginal
$\min\{F_\chi(y):y_i=v\}$, with $G_i$ replaced by $\{v\}$.
\end{lemma}

The proof (Appendix~\ref{app:recourse-cuts}) encodes each coordinate by
threshold variables and applies Proposition~\ref{prop:cr-cut}(a); arcs of
infinite capacity exclude the binary vectors that encode no grid node.
Threshold encodings, and the reduction of pairwise problems that are
submodular with respect to label orders to one minimum cut, are classical
\cite[Sections~3--4]{SchlesingerFlach2006}, \cite{Ishikawa2003}.
Lemma~\ref{lem:chaincut} is this construction for quadratic pair terms; what
matters here is that the unary terms $\chi_i$ are arbitrary, so that they
can be the negative corrections $-d_i$ on nonuniform grids.

\begin{theorem}[Balanced box quadratics without a width parameter]\label{thm:balanced}
Let $F$ be a balanced rational quadratic on a mixed box $X$ with input
length $I$. Run CT and EX (Algorithm~\ref{alg:ex}) with $\beta$, a
minimizer and all min-marginals computed by Lemma~\ref{lem:chaincut} instead
of Lemma~\ref{lem:dp}.
\begin{enumerate}[label=(\alph*)]
\item On every instance, CT$(2^{-q})$ halts with a feasible rational point
and a valid path certificate of gap at most $2^{-q}$, in which the dynamic
programs are replaced by minimum cuts with flows of equal value; and EX
halts with a global minimizer and $\OPT$.
\item If $F$ has point growth with condition number $\kappa$, then
CT$(2^{-q})$ uses $(n\kappa(I+q+1))^{O(1)}$ bit operations and EX uses
$(n\kappa(I+1))^{O(1)}$ bit operations.
\end{enumerate}
No decomposition is used, $g$ and $\kappa$ are not inputs, and the validity
of every certificate uses neither growth nor uniqueness.
\end{theorem}

\begin{proof}
(a) In CT and EX, and in the proofs of Propositions~\ref{prop:cellwise}
and~\ref{prop:filter}, Theorem~\ref{thm:certificate},
Lemmas~\ref{lem:inv} and~\ref{lem:states} and
Theorems~\ref{thm:approx}, \ref{thm:transfer} and~\ref{thm:exact}, the tree
decomposition enters only through Lemma~\ref{lem:dp}, which computes
$\beta$, a minimizer of $Q$ and the min-marginals $m_i(v)$ on the current
grid (if $P=\emptyset$, CT minimizes $F$ over the grid
$\prod_i\{\ell_i,u_i\}$, on which $Q=F$). Everything else uses the data of
$F$ directly: the condition of Definition~\ref{def:curvature} holds with
$L_i=\max\{H_{ii},0\}$, because the restriction of $F$ to a coordinate line
has second derivative $H_{ii}$, and REC (Algorithm~\ref{alg:rec}) and
Proposition~\ref{prop:accept} work with $F$ and its gradient. Since $D$ is a
sum of unary terms, $Q=F_\chi$ with $\chi_i=-d_i$, so
Lemma~\ref{lem:chaincut} computes $\beta$ and a minimizer with one cut and
each $m_i(v)$ with one further cut. A flow of equal value certifies each of
these values from below by weak duality, which is what conditions (C1)
and~(C2) of a path certificate require. Hence the validity and termination
statements of these results hold verbatim, including the abort rule of CT
and, under weighted growth, the success of a trial with
$2^\mu\le6\sqrt{\bar\kappa}$; only the operation counts change.

(b) It remains to count cut computations and bit lengths in place of table
operations; the count is in Appendix~\ref{app:recourse-cuts}.
\end{proof}

\begin{corollary}[Balanced after deleting a supplied core]\label{cor:core}
Let $F$ and $X$ be as in Theorem~\ref{thm:balanced}, except that $F$ need not
be balanced. Suppose a set $\mathcal K$ of $k$ coordinates is supplied such
that $o_io_jH_{ij}\le0$ for some signs $o$ and all $i\ne j$ outside
$\mathcal K$. Then Theorem~\ref{thm:balanced} holds for $F$, with the work
bounds of part~(b) multiplied by $K^k$, where $K=O(\sqrt\kappa\log n)$
bounds the nodes per coordinate in every trial that CT runs.
\end{corollary}

\begin{proof}
Fixing the core coordinates at grid nodes changes only unary and constant
terms of the remaining quadratic, which is balanced. Hence $\beta$, a
minimizer and each min-marginal are minima, over at most $K^k$ vectors of
core nodes, of problems solved by Lemma~\ref{lem:chaincut}; a min-marginal
of a core coordinate fixes that coordinate. Every trial that CT runs has at most
$K_\mu\le K_{\mu^*}$ nodes per coordinate, where
$2^{\mu^*}\le6\sqrt{\bar\kappa}\le6\sqrt\kappa$ (proof of
Theorem~\ref{thm:approx}(b) and Appendix~\ref{app:recourse-cuts}), so
$K=K_{\mu^*}=O(\sqrt\kappa\log n)$.
The rest of the proof of Theorem~\ref{thm:balanced} is unchanged.
\end{proof}

Corollary~\ref{cor:core} complements Section~\ref{sec:cuts}: residual
coordinates may now have positive diagonal, but they are gridded and filtered
like the core, so their curvature enters $\kappa$, and point growth of $F$
is needed instead of growth of $V$ in the core coordinates; the residual
graph may still be dense.
Section~\ref{sec:related} compares Theorem~\ref{thm:balanced} with methods
for continuous submodular minimization.

